\section{Conclusion}
We presented CRF-SAM3, a Ground--Recall--Verify framework for referring video object segmentation. Its central premise is that object discovery and referent identification require different semantic granularity. An expression-conditioned target-query grounder uses mask supervision only during training to form a referent-aware open-vocabulary recall concept from RGB frames and language. SAM3 then exposes multiple temporally coherent candidate tubes, while a candidate-aware Qwen3-VL verifier preserves the complete expression for final discrimination through listwise competition, hard-negative ranking, and mask-aware feature alignment.

Across MeViS, Ref-DAVIS17, Ref-Youtube-VOS, and ReVOS, the framework shows that explicit candidate construction provides both competitive segmentation performance and a useful diagnostic interface. The MeViS result is particularly encouraging for motion-centric settings, while the complete ReVOS validation run reaches 59.74 $\mathcal{J}\&\mathcal{F}$ and exposes a clear remaining gap on reasoning-oriented expressions. Candidate-space analysis further shows that prompt-sensitive recall creates an upper bound for later reasoning, whereas verifier training substantially improves discrimination when a valid target tube is already present. These observations motivate future work that couples concept formation more directly to downstream candidate quality while retaining the interpretability of separate recall and verification stages.
